\section{Retained local, endpoint and pointwise statements}
\label{app:v57-retained-front}


\paragraph{The original record and its density.}
Let $T_R$ be the physical collision map and $F_R^*$ its actual first
return to $Y_R^*$. The collision probability is $\nu$ and the return
probability is $\nu_R^*=c^{-1}\nu|_{Y_R^*}$, where
$c=91/(10000\pi)$. For $n$ genuine returns write
$J_{n,R}=(K_{n,R},N_{n,R},T_{n,R})$. Define $p_{n,R}$ by
\[
 \nu_R^*\{K_{n,R}=k,N_{n,R}=m,T_{n,R}\in du\}
                  =p_{n,R}(k,m,u)\,du.
\]
Its existence on the complete, unprotected source is
Theorem~\ref{thm:v44-complete-density-bound}. The counting coordinates
are $k\in\Z^2$ and $m\ge n$; only the roof coordinate $u$ is continuous.

\begin{leadtheorem}[An explicit weak endpoint for the original arithmetic density]
\label{thm:intro-v49-local-TV}
Let $\mathcal L_{m,R}$ be the evaluated finite transition kernel in
\eqref{eq:v41-transition-kernel}. For every fixed $h>0$,
\[
 \lim_{m\to\infty}\sup_{R,\,1\le n\le m,\,k,t}
 \frac1h\int_t^{t+h}
 |m^2p_{n,R}(k,m,u)-\mathcal L_{m,R}(k_1,k_2,u,n)|\,du=0.
\]
The original exact return index is retained. At fixed radius and on
central compact sets the main term is
$c\mathfrak a_R(k,n,m)g_{\Omega_R}$ in the stated collision
normalization. There is no zero-residue assumption.

Put $p_c=145/144$. The normalized original density $P=m^2p_{n,R}$
satisfies
\[
 \sup_{R,n,k,m}\sup_{L>0} L^{p_c}
    |\{u\in[t,t+h]:P(u)>L\}|\le C(1+h).
\]
For every $1<q<p_c$,
\[
 \sup_{R,n,k,t}\frac1h\int_t^{t+h}
 |m^2p_{n,R}(k,m,u)-\mathcal L_{m,R}(k_1,k_2,u,n)|^q\,du\to0.
\]
For $\beta>1$, the same assertion holds with the $q$th power replaced
by the convex function $\Phi_\beta$ in
\eqref{eq:v55-young-function}, which is comparable at infinity to
$x^{p_c}/[\log(e+x)]^\beta$. This is a weak endpoint and an
endpoint-modular law, not a strong $L^{p_c}$ assertion.

On fixed windows with pointwise arithmetic reference floor
$\mathcal L_{m,R}\ge d>0$, the unchanged conditional roof likelihood
converges to one in every $L^q$ of the reference law, $1<q<p_c$,
and in each of these endpoint modulars. Forward relative entropy
and forward R\'enyi divergences of orders in $(1,p_c)$ tend to zero.
No collision-count rate or forward essential likelihood estimate
is claimed.

There is also a resolution-uniform localization. For each large fixed
$B$, let $O_{B,m}$ be the positive physical maximal set of
\eqref{eq:v56-band-threshold}. Its length on the roof line is at most
$C(1+h)B^{-1/384}$, and its source mass on the displayed window has
collision limsup at most the same bound. At every anchor outside
$O_{B,m}$, all subinterval averages of the scalar and collision-path
raw errors are bounded by $r_{B,m}$, with
$\limsup_m r_{B,m}\le CB^{-1/384}$. The subintervals have no minimum
length. The same-roof statement holds almost everywhere on that
complement; it does not assert height control inside $O_{B,m}$.

The same limit is uniform over all measurable roof selectors of
modulus at most one supported in the displayed interval. On bounded
central regions at a fixed return index, it gives total-variation
approximation of the complete mixed measure. If the positive-part
transition mass of a roof interval is at least $dh$, $d>0$, the exact
conditional roof law is approximated in total variation by its
normalized positive-part transition density.
\end{leadtheorem}
\begin{proof}
Apply Theorem~\ref{thm:v49-raw-local-TV} and
Corollaries~\ref{cor:v49-roof-selector-TV}, \ref{cor:v49-mixed-TV}
and~\ref{cor:v49-conditional-roof-TV}, followed by
Theorems~\ref{thm:v54-uniform-integrability}, \ref{thm:v54-raw-Lq}
and~\ref{thm:v54-forward-likelihood}, followed by
Corollary~\ref{cor:v55-physical-endpoint} and
Theorems~\ref{thm:v55-raw-endpoint} and
\ref{thm:v55-endpoint-likelihood}, together with Corollary~\ref{cor:v56-ordered-localization} and Theorem~\ref{thm:v56-all-resolutions}.
\end{proof}


\begin{leadtheorem}[Pointwise arithmetic lower law and positive concentration]
\label{thm:intro-v51-lower}
Write $P(u)=m^2p_{n,R}(k,m,u)$ and
$G(u)=\mathcal L_{m,R}(k_1,k_2,u,n)$. Then
\[
 \sup_{R,1\le n\le m,k}\operatorname*{ess\,sup}_u(G(u)-P(u))_+
                                                       \longrightarrow0.
\]
More precisely, at each sufficiently large fixed band $B$, let
$\mathfrak b_B=m^2b^{A_0B^{-1/12},1}_{n,k,m,R}\ge0$ be the
original positive physical remainder. Then
\[
 \limsup_{m\to\infty}\sup_{R,n,k}
                  \|P-G-\mathfrak b_B\|_\infty\le CB^{-1/192}.
\]
The full two-sided pointwise theorem on central compact sets is
equivalent to the vanishing of the positive height of $\mathfrak b_B$
in the ordered limit $m\to\infty$, then $B\to\infty$.
This condition is also equivalent to the separate vanishing of its
positive incidence and clearance components.

Where the transition kernel is bounded below by $d>0$, the original
roof density is eventually at least $d/(2m^2)$ almost everywhere.
On fixed roof windows with that lower bound, the exact conditional
roof law dominates $(1-o(1))$ times the normalized arithmetic reference
as a positive measure, uniformly in all targets.
\end{leadtheorem}
\begin{proof}
Theorem~\ref{thm:v51-positive-error},
Corollaries~\ref{cor:v51-positive-criterion} and
\ref{cor:v51-lower-denominators}, and
Theorem~\ref{thm:v51-conditional-minorization} give the assertions.
\end{proof}

\begin{leadtheorem}[Path-valued raw limits and same-roof conditional bridges]
\label{thm:intro-v52-path}
Let $\mathsf Q_{m,R}^{n,k,u}$ be the conditional law of the pinned
collision path given the original exact discrete labels and its own
roof value $u$, and let $\mathsf W_R$ be the Gaussian bridge law with
covariance $\Omega_R$. The norm $\mathrm{BL}^*$ is the bounded-Lipschitz
dual norm on continuous path space. For every fixed $h>0$,
\[
 \sup_{R,n,k,t}\frac1h\int_t^{t+h}
 \left\|m^2p_{n,R}(k,m,u)\mathsf Q_{m,R}^{n,k,u}
      -\mathcal L_{m,R}(k_1,k_2,u,n)\mathsf W_R\right\|_{\mathrm{BL}^*}
 \,du\longrightarrow0.
\]
At every sufficiently large fixed band $B$, the complete path-valued
error differs, in essential supremum of this norm, from the original
positive physical path measure by at most $CB^{-1/192}$ after the
collision limsup. The same physical measure works simultaneously for
all unit bounded-Lipschitz path tests.

On fixed central target sets, the displayed local limit also holds
for the actual return path and the Gaussian bridge of covariance
$D_R$. Whenever the positive-part arithmetic reference mass on the
roof interval is at least $dh$, $d>0$, the conditional bridge kernel
at $u$ converges in mean bounded-Lipschitz distance, uniformly, under
the original conditional roof law. In particular it converges uniformly
outside sets of roof values whose conditional probabilities tend to
zero. These kernels fix the same exact event and the same roof value;
no pointwise conditioning event is replaced by a window event.

For each fixed central window and some $a>0$, the unnormalized exact
source also satisfies
\[
 \frac{m^2}{h}\int_E
 e^{a(\|\mathcal B_{m,R}\|_\infty+\|\mathcal Y_{n,R}\|_\infty)}
                                                 d\nu_R^*\le C.
\]
Consequently the path-valued local laws hold for polynomial-growth
Lipschitz tests, and for every finite $p\ge1$,
\[
 \int W_p(\mathsf Q^u,\mathsf W_R)^p\,d\pi_m(u)
  +\int W_p(\mathsf R^u,\mathsf V_R)^p\,d\pi_m(u)\longrightarrow0.
\]
Here $\pi_m$ is the unchanged conditional roof law and $W_p$ uses the
unbounded path supremum metric. Conditional covariance kernels and
path-maximum moments converge in the corresponding roof mean. The
same unbounded-cost assertions under the reference roof law additionally
require a pointwise positive reference floor on the window.

The scalar positive-height criterion of Theorem~\ref{thm:intro-v51-lower}
also suffices for the uniform essential-supremum same-roof bridge on
positive-reference central sets. This implication does not assert that
the remaining height criterion has been proved.
\end{leadtheorem}
\begin{proof}
Theorem~\ref{thm:v52-path-remainder}, Lemma~\ref{lem:v52-return-transfer},
Theorem~\ref{thm:v52-same-roof-mean},
Corollary~\ref{cor:v52-full-roof-sets} and
Proposition~\ref{prop:v52-pointwise-path-charge} give the bounded-test
assertions. Theorem~\ref{thm:v53-exponential-maximal},
Theorem~\ref{thm:v53-unbounded-raw-paths},
Theorem~\ref{thm:v53-same-roof-Wp} and
Corollary~\ref{cor:v53-conditional-moments} prove the strengthened
moment and transport conclusions.
\end{proof}

\paragraph{A direct route to the path numerator.}
The new route requires four interfaces, not a new inducing scheme.
The endpoint influence estimate in
Lemma~\ref{lem:v46-regular-endpoints} bounds the whole path map in
Lemma~\ref{lem:v52-path-endpoint-budget}. The chronological peak
calculation and unnormalized fourth moments in
Section~\ref{sec:v41-bridges} yield fixed-band inversion uniformly over
path tests. The protected correction, all-depth decision estimate and
physical local variation then prove
Theorem~\ref{thm:v52-path-remainder}. Section~\ref{sec:v52-same-roof}
disintegrates this actual numerator and transfers the pinned clock.
The norm integrates an absolute path-dual error in the roof coordinate;
it is neither path-space total variation nor a two-sided scalar
essential-height theorem.

\paragraph{A cap-free route to the explicit endpoint.}
The new scalar route uses only three finite inputs: the local all-margin
mass bound of Lemma~\ref{lem:v54-fixed-band-windows}, the polynomial
protected height of Proposition~\ref{prop:v54-protected-height}, and
the exact global strip mass of Lemma~\ref{lem:v55-global-layer}.
At large $m$ relative to $\log(1/\varepsilon)$, the local estimate
absorbs its exponential error. At smaller $m$, the global mass is
stronger. Lemma~\ref{lem:v55-two-regimes} therefore gives
$C\varepsilon^{1/16}$ with no remaining finite-count error.
The positive-source split then gives a density-weighted tail
$CL^{-1/144}$ and a weak $L^{145/144}$ estimate. The full exponential
height cap is no longer an input to this route, and the exponent does
not depend on a ratio of spectral decay to geometric height growth.
All prior proofs, including the earlier cap-based argument, are retained.

Section~\ref{sec:v55-endpoint-laws} proves, for each physical component,
\[
 \sup_{R,n,k,t,m}\int_t^{t+h}
 |m^2b^{\varepsilon,j,w}_{n,k,m,R}(u)|^q\,du
 \le C_qM^q(1+h)\varepsilon^{9(p_c-q)},\quad 1<q<p_c.
\]
The same estimate holds for its complete smoothing correction at every
bandwidth. These finite inequalities allow a count-dependent protection
scale, but do not use a count-dependent spectral band. The scalar and
path-valued endpoint laws and the original roof likelihood follow from
the same distribution-tail argument. The arithmetic factor and finite
transition kernel are permanent parts of the statements, not terms
assumed to cancel.

\paragraph{From local variation to a one-sided pointwise theorem.}
The new argument uses the actual positive source partition, not the
invalid implication from local $L^1$ convergence to uniform convergence.
Section~\ref{sec:v51-positive-remainder} first estimates the convolution
of the physical remainder in essential supremum, with the reconstruction
band fixed. The protected and complete section-decision corrections
already have pointwise bounds. Subtracting these quantities leaves a
positive physical source, which proves the lower law and the sharpened
positive-height criterion. Section~\ref{sec:v51-roof-minorization}
then obtains exact conditional minorization. Positive spikes and the zero-residue specialization do not follow
from those statements. The new path numerator is proved separately in
Section~\ref{sec:v52-path-inversion}; its same-roof convergence in mean
is distinguished from the uniform pointwise conclusion which still
requires the scalar height criterion.

\paragraph{The physical remainder and the ordered estimate.}
For a locally integrable density set
$\|f\|_{1,h}^{\rm loc}=\sup_t\int_t^{t+h}|f|$.
The actual physical remainder $b^{\varepsilon,w}$ contains all
histories with an incidence or clearance defect, whether or not they
also carry section-decision defects. For arbitrary $|w|\le M$,
Theorem~\ref{thm:v49-physical-remainder-local} proves
\[
 \sup_{R,n,k}m^2\|b^{\varepsilon,w}_{n,k,m,R}\|_{1,h}^{\rm loc}
 \le CM\varepsilon^{1/16}(h+B_*^{-1})
       +C_{B_*}M(1+B_*h)m^3\rho_{B_*}^{m/2}.
\]
The auxiliary band $B_*$ is fixed, but the defect width and mark may
vary with $m$. The full correction, with no guard on the original
density, therefore satisfies
\[
 \limsup_{m\to\infty}\sup_{R,n,k}
 \frac{m^2}{h}\|p_{n,R}(k,m,\cdot)-K_B*p_{n,R}(k,m,\cdot)\|_{1,h}^{\rm loc}
                         \le CB^{-1/192}.
\]
Here $B$ is fixed before the collision limsup and
$\varepsilon(B)=AB^{-1/12}$. The auxiliary band is enlarged only
after that limsup. No rate in $m$ or growing-band spectral estimate
is asserted.

\paragraph{A short proof route.}
Section~\ref{sec:v49-physical-layers} proves the new geometric
multiplier: clearance is tested at the next collision, where its
boundaries have positive slope and are transverse to stable curves.
It combines this with the previously proved moving-peak interpolation.
Section~\ref{sec:v49-physical-variation} sums all physical marks
with a uniform finite-count remainder and estimates their signed
correction in the local density norm.
Section~\ref{sec:v49-raw-local-TV} combines this with the already
proved protected and section-decision corrections and the fixed-band
arithmetic kernel. The earlier leading statements are compiled in
Appendix~\ref{app:v49-retained-leading}; all earlier core proofs remain.

\paragraph{Pointwise raw inversion remains the original target.}
Local total variation is stronger than testing a fixed roof interval:
the test may be chosen measurably after the collision count. It is not
an essential-supremum statement. The same original pointwise criterion
\eqref{eq:v48-physical-residual-criterion} still requires control of
possible narrow physical-boundary spikes. We have neither changed that
criterion nor replaced its arithmetic main term. The scalar conditional result of Theorem~\ref{thm:intro-v49-local-TV}
is for the unchanged event with $T_{n,R}$ in a fixed interval.
Theorem~\ref{thm:intro-v52-path} estimates the disintegrated path
kernels at their own roof values in conditional mean; it does not
assert convergence at every prescribed roof value.

\paragraph{Relation to previous local limit methods.}
The collision-space framework and piecewise multiplier criterion are
those of \cite{DZ11,DZ,DPZ}; the comparisons with collision and
suspension local limits remain in
Section~\ref{sec:v35-literature-routes}. The new conclusion is not
obtained by differentiating their interval statements. It uses the
physical thin-layer estimate and previously proved protected
regularity to control an absolute density norm on the full source.
No claim of priority for general mixing local limits is made.

